\section{An outward coupling conditioned on two arms}
\label{sec:coupling}

The electric correlations will compare a pair of points at a growing
separation with the two sites adjacent to a single interface segment.
The comparison requires more than an estimate for an arm probability:
after conditioning the two colors to reach a distant rim, the exterior
configurations must become identical with high probability.  We prove
this statement here.  The same coupling will also compare a segment
with a small disk.

All distances in this section are in tile units.  A cut is the boundary
of a finite union of whole tiles.  A strand is cut at every port on this
boundary.  We use square cuts for the successive coupling trials and
allow either a square or a round cut for the final target.  Round cuts
are simple whole-tile approximations to circles, at bounded Hausdorff
distance from those circles.  In particular, reaching a cut means
reaching an exposed gap of the appropriate color.  Boundary
identifications do not supply edges of a color path.

\subsection{Three starting configurations}

Fix a center and a target cut of radius $R$.  We consider the following
starts, contained well inside that cut.
\begin{enumerate}[label=\textup{(\roman*)}]
\item A tile side $e$.  Its two end sites have opposite colors, and the
associated full medial segment crosses that side.  The event $S_R(e)$
requires an actual path of each color from its end site to the target
cut.
\item Two sites $u,v$ of opposite colors, with $|u-v|\asymp r$ and both
within distance $O(r)$ of the center.  The event $S_R(u,v)$ requires
their respective color paths to reach the target cut.
\item A round whole-tile disk of radius $r$.  The event $S_R(r)$ requires
both colors to cross the intervening annulus, starting at gaps of the
disk boundary.
\end{enumerate}
We sometimes write $S_R$ when the start is specified separately.  Every
path in this definition is restricted to the region between the start
and its first arrival at the target cut.

These definitions have useful descriptions in the strand picture.
For a point start, its color region reaches the target cut if and only
if no completed loop inside that cut surrounds the point.  Indeed, a
region that does not reach the cut has a boundary loop separating it
from the cut; conversely, an actual color path cannot cross a loop.
Thus in cases \textup{(i)} and \textup{(ii)}, $S_R$ says that no
completed loop surrounds either starting site.  In case \textup{(iii)},
$S_R(r)$ is equivalent to the existence of a strand subarc joining the
two rims.  To see the latter assertion, cut all strands at both rims.
An arc with both ends on the same rim cuts off a region incident only
to that rim.  If there are no through-arcs, at most one color can have
a region incident to both rims.  If there are through-arcs, the regions
between consecutive through-arcs give paths of both colors between
the rims.  The rounded drawing makes this a statement about disjoint
arcs, even when the straight medial drawing has a double visit.

Write $\pi_1(r)$ for the plane one-arm probability to radius $r$, and
$\pi_2(r,R)$ for the plane alternating two-arm probability across a
square annulus.  Their estimates from Lemmas~\ref{lem:one-arm}
and~\ref{lem:two-arm-log} are available below.  A fixed change of a
radius, or of the shape of a rim between the round and square choices,
changes the following comparisons by a constant only.

\begin{lemma}[Probabilities of the three starts]
\label{lem:start-probabilities}
There is a fixed minimum radius $r_0$ such that, for $R$ sufficiently
larger than the start,
\begin{equation}\label{eq:start-probabilities}
 \Pp(S_R)\asymp
 \begin{cases}
   \pi_2(r_0,R),&\text{for a tile-side start},\\
   \pi_1(r)^2\pi_2(r,R),&\text{for two sites at separation }\asymp r,\\
   \pi_2(r,R),&\text{for a disk of radius }r.
 \end{cases}
\end{equation}
The same constants apply in an ambient FK law with a full bulk collar
beyond the target cut, including after a nonnegative integrable
reweighting supported beyond that collar.  In the second line, the
constants can depend on the fixed comparability bounds for the two
sites.
\end{lemma}

\begin{proof}
For the upper bound in the second line, choose disjoint boxes of radius
$c r$ about the two sites, with $c>0$ small enough.  Each site must have
a one-arm connection to its own box boundary.  Outside a square of
radius $C r$ about the center there must also be an alternating pair
of arms extending to a fixed fraction of $R$.  Leave fixed-ratio
collars between these three observations.  Bounded-density comparison
through these collars, as recorded in Input~\ref{inp:fk}, bounds their
joint probability by
$C\pi_1(c r)^2\pi_2(C r,R/C)$.  One-arm extension and
quasimultiplicativity give the asserted upper bound.

For the lower bound, require a one-arm event in each small box with a
landing extension pointing into a designated corridor.  Require
alternating arms in the outer annulus with separated inner landings.
The two corridors can be chosen disjoint, one for each color.  Uniform
crossing estimates and same-color FKG connect the corresponding
landings at a positive fixed cost.  The collars give a lower comparison
for the events before this gluing, and quasimultiplicativity removes
the fixed changes of radii.  This proves the second line.  It also
explains why no favorable boundary partition is required: each
crossing estimate is uniform in the partition induced by the already
observed edges.

For a tile-side start, the initial part of the two paths can be imposed
inside a fixed box by finite energy, after which the same two-arm
extension argument applies.  Cropping gives the reverse inequality.
For a disk start, crop or extend at its two rims and use the annular
two-arm event directly.  These arguments take place before the
remaining bulk collar.  Conditioning on all data beyond that collar
proves the bounds uniformly in the exterior state.  Integration proves
the assertion for a nonnegative exterior reweighting.
\end{proof}

\begin{proposition}[Outward conditioned coupling]
\label{prop:outward}
Fix $q\in[1,4)$.  Consider one ambient FK law whose graph is the full
tile lattice throughout the region bounded by the target cut and in
a fixed-ratio bulk collar beyond it.  All imposed wires and
nontrivial identifications are supported outside this region and
collar.  The law may be multiplied by a common
nonnegative integrable function supported beyond the collar, with
positive finite normalizing expectation.  Let two starts of the types
above be contained in a disk of radius $s$, where $1\le s\le R/C$.
Their laws, conditioned on their respective events $S_R$, have a
coupling such that, except with probability
\begin{equation}\label{eq:outward-error}
 C(s/R)^c,
\end{equation}
the configurations agree on and outside an annular band lying between
the starts and the target cut.  There are exactly two strand subarcs
joining the rims of this band.  Each start reaches the same
transmitted primal region and the same transmitted dual region of
the band.  Consequently the exterior sees the same pairing of its
exposed strand ends in the two configurations.

Here $c>0$ and $C<\infty$ depend only on $q$, the fixed collar
geometry, and the fixed comparability constants of the starts.  The
assertion is uniform in the exterior reweighting.
\end{proposition}

The proof constructs an annular band which has positive probability
and transmits exactly one region of each color.  We first explain its
geometry and the gluing that preserves an arbitrary event on its
central edges.  We then show that a conditioned exploration has enough
opportunities to use this band.

\subsection{A band with exactly two through-arcs}

\begin{lemma}[The reference band]
\label{lem:reference-band}
At every sufficiently large scale $l$, there is a fixed-shape square
annular band $\mathcal B_l$, of inner and outer radii comparable to
$l$, and an event of probability bounded below in the plane FK law
with the following properties.  Both colors cross $\mathcal B_l$;
there are exactly two strand through-arcs; and each of the two
crossing color regions has separated landing extensions beyond both
rims.  The extension boxes can be chosen a positive relative distance
from the central band and from each other.  The constants and the
relative sizes of all these regions are independent of $l$.
\end{lemma}

\begin{proof}
Choose one radial lane on the top flat sides of two concentric squares
and another on their bottom flat sides.  All lane widths and the gaps
between lanes are fixed positive multiples of $l$.  The top lane will
carry primal crossings and the bottom lane dual crossings.  In each
lane choose three parallel rectangular strips: a main strip and a
flanking strip on each side.  Require longitudinal crossings of all
three strips.  The main crossing is extended beyond both square cuts
into prescribed, separated landing boxes.  Inside the central annulus,
require transverse same-color crossings joining each flanking
crossing to the main crossing.  A transverse crossing runs across the
full width of a flank strip, so it meets every longitudinal crossing
of that strip.  Crossing estimates and FKG give a positive fixed
probability for these requirements.  The regions for the two colors
are disjoint, so the colors can be treated successively using the
conditional crossing bounds.

In each flanking strip choose the extremal longitudinal crossing
closest to the main strip and expose it from the main-strip side.
The transverse joining paths can be stopped at their first contact
with this crossing.  Thus the resulting event still contains the
crossing event just imposed, but its verification does not inspect
the two corridors between opposing flanks of different colors.
This is the usual extremal-crossing exposure: the path and the
configuration on its searched side determine the extremal path,
whereas edges strictly on its other side retain an FK conditional
law.  The main roads and their tethers are connected within the
central annulus.

We next join the two color lanes around both sides of the annulus.
In either remaining corridor, the dual flank is one target side.
The other target side consists of dual cells adjacent to the
corridor side of the primal flank.  We give the lattice detail needed
for this second side, since replacing it by a wired identification
would not produce the required barrier.

Take the central square cuts along dual-lattice rim lines, adjusting
by boundedly many rows.  The primal flank starts half a mesh before
the first rim line and ends half a mesh after the second.  Retain its
portion ending at its first visit to the ending line and beginning
at its last visit to the starting line before that ending-line visit.
Its interior then avoids both lines.  Follow
the centers of the dual cells incident to the path on its corridor
side.  At a turn, continue through the centers in that side's sector
around the primal vertex.  Consecutive centers are nearest neighbors;
none of these dual steps crosses an edge of the primal path.  The
resulting walk reaches both dual rim lines.  Retain a crossing
subwalk with its interior off those lines and erase loops.  This
gives a simple dual side arc, still in the flank strip and incident
to the primal path.  All its edges are on the unsearched side of
the extremal exposure.  Together with the dual flank and the two
rim intervals, it bounds a dual-lattice quadrilateral in the
unrevealed corridor.

The extremal distance between the two target sides of this
quadrilateral is uniformly bounded above.  Indeed, a family of
square contours at intermediate radii crosses the fixed-width gap
between the lane strips.  Continue each contour to its first hits
of the two target sides.  There are order $l$ such disjoint
contours, each of length $O(l)$.  The corresponding parallel paths,
or the extremal-length variational formula, give a constant upper
bound independent of the jagged flank geometry.  The quadrilateral
crossing theorem in Input~\ref{inp:fk} therefore gives an actual dual
crossing from the dual flank to a dual cell incident to the primal
flank, with conditional probability bounded below.  Already open
edges on the dual flank only help.  The two corridors are treated
in succession, each with its induced boundary partition.

The two corridor crossings, the tethers, and the two main roads now
form a closed route around the inner square.  One portion is an
actual primal path and the other an actual dual path.  They are
joined by exactly two steps between adjacent sites of opposite
colors.  Its winding number about the inner square is one: its
portions lie successively in the top lane, the right corridor, the
bottom lane, and the left corridor.  Erasing detours within any of
these four regions does not change that winding number.  Such a
route is a barrier even if it was not initially chosen simple.

Figure~\ref{fig:barrier} shows this closed route and the two
adjacency steps at which it can meet medial strands.
An actual color path does not cross a medial strand.  Each of the
two steps between adjacent opposite-color sites crosses exactly
one medial segment.  Every strand through-arc of a slightly wider
central band must intersect the closed route, because the route
separates the two rims.  Distinct through-arcs are disjoint in the
rounded drawing, and therefore use different crossing segments of
the route.  There are at most two through-arcs.  The longitudinal
roads give crossings of both colors, so the elementary annular
description above gives at least two.  There are exactly two.
Enlarge the central band by fixed margins while keeping both ends
of the main roads and their landing boxes outside it.  This gives
the stated extensions without changing the barrier argument.
Every construction has used finitely many crossing requirements
of fixed geometry, so their total probability has a positive
scale-independent lower bound.
\end{proof}

\begin{figure}[tbp]
\centering
\begin{tikzpicture}[x=0.86cm,y=0.86cm,>=Latex,
  every node/.style={font=\small}]
  \fill[black!5] (-3,-2) rectangle (3,2);
  \fill[white] (-1,-.65) rectangle (1,.65);
  \draw[black!55] (-3,-2) rectangle (3,2);
  \draw[black!55] (-1,-.65) rectangle (1,.65);
  \node at (0,0) {inner rim};
  \node[anchor=west] at (3.12,-1.8) {outer rim};
  \draw[black!55] (3.03,-1.8)--(3.1,-1.8);

  \draw[blue!70!black,line width=1.35pt]
    (0,.65)--(0,2.35);
  \draw[red!75!black,line width=1.35pt,densely dashed]
    (0,-.65)--(0,-2.35);
  \node[anchor=west,text=blue!70!black] at (.12,2.27)
    {primal road};
  \node[anchor=west,text=red!75!black] at (.12,-2.27)
    {dual road};

  \draw[blue!70!black,line width=1.6pt]
    (-1.62,1.00)--(-.8,1.22)--(0,1.25)--(.8,1.22)--(1.62,1.00);
  \draw[red!75!black,line width=1.6pt,densely dashed]
    (1.66,.95)--(2,.15)--(2,-1.2)--(-2,-1.2)--(-2,.15)--(-1.66,.95);
  \draw[black,line width=1.1pt] (1.62,1.00)--(1.66,.95);
  \draw[black,line width=1.1pt] (-1.62,1.00)--(-1.66,.95);

  \draw[black,line width=.8pt]
    (1,.45).. controls (1.26,.67) and (1.40,.79) .. (1.64,.975)
      .. controls (1.98,1.24) and (2.52,1.66) .. (2.9,2);
  \draw[black,line width=.8pt]
    (-1,.45).. controls (-1.26,.67) and (-1.40,.79) .. (-1.64,.975)
      .. controls (-1.98,1.24) and (-2.52,1.66) .. (-2.9,2);
  \draw[black,fill=white,line width=.6pt] (1.64,.975) circle (.072);
  \draw[black,fill=white,line width=.6pt] (-1.64,.975) circle (.072);
  \node[anchor=south east] at (1.59,1.10) {$a_2$};
  \node[anchor=south west] at (-1.59,1.10) {$a_1$};
  \node[anchor=west] at (3.12,.55) {through-arcs};
  \draw[black!65,->] (3.05,.64)--(2.27,1.46);
  \node at (1.5,-1.65) {barrier};
\end{tikzpicture}

\caption{The topological consequence of the reference-band
construction.  A primal path and a dual path form a route of winding
number one around the inner rim, with exactly two adjacency steps
$a_1,a_2$ between colors.  Only those steps can cross medial strands.
Hence at most two strand arcs can cross the band; the two radial
color roads force both to be present.  The picture is schematic:
the adjacency steps are enlarged, the flanking crossings used to
construct the route are omitted, and neither the color paths nor
the medial arcs are asserted to be straight or smooth.}
\label{fig:barrier}
\end{figure}

\subsection{Landing extensions and gluing}

We use a strengthened form of arm separation which retains the
crossing clusters.  For a fixed-ratio bulk annulus and every
$\varepsilon>0$, the probability in the plane FK law that some crossing arms admit no
separated landing extensions in their respective original color
clusters is at most $\varepsilon$, when the separation parameter is
chosen small enough.  This is precisely the event $B_n$ preceding
Lemma~5.9 of \cite{CDCH}.  The replacement arms there are required to
belong to the same respective primal or dual clusters.  The proof
uses the quadrilateral crossing bounds, including at the exposed
jagged sides, the small- and large-extremal-length bounds, and an
upper bound on the number of crossing components.  These are all
available for $1\le q<4$ by \cite[Theorem~1.2 and Section~6]{DMT};
indeed its proof of Proposition~6.2 invokes this argument.  Thus we
use the proof at its stated probabilistic inputs, rather than an
Ising-specific identity.

We need two geometric features of these extensions.  Their endpoints
stay away from the corners of the square cut, and the extension
boxes can be much smaller than the separation between different
landings.  In the proof of \cite[Lemma~5.9]{CDCH}, separation occurs
at the scale denoted there by $\delta^3 n$, whereas extensions are
constructed at smaller scales $\delta^k n$, $k\ge4$.  Only finitely
many such scales are needed for a specified error tolerance.  The
arms consequently protrude a positive relative distance beyond
their landing cut, inside disjoint boxes having a fixed-factor
margin between them.  All constants may depend on the error
tolerance.  Bounded lattice rounding is harmless once the trial
scale exceeds a corresponding fixed minimum.

Suppose a past has been exposed up to a square cut of radius $l$,
and the conditional, unconditioned-FK probability that the specified
start reaches a larger cut of radius $8l$ is at least $h>0$.
Here and below the factors $8$ and $16$ may be increased by fixed
amounts to accommodate collars.  Apply the preceding landing
statement on an annulus strictly between those two cuts.  Its
failure event can be assigned ordinary plane probability less than
$h/(2C)$, where $C$ is the collar-comparison constant.  Its
conditional probability given the past is then less than $h/2$.
Thus, with conditional probability at least $h/2$, the two
start-connected clusters have separated outer extensions before
radius $16l$.  The crossing portions are taken after their last
visits to the annulus's inner rim.  Replacing those portions within
their same clusters preserves connection to the start.  A finite
cover of the landing rim by small boxes allows the landing
locations to be restricted to one pair of boxes, at a cost depending
on $h$ only.

There is a corresponding construction facing inward from the
remaining annulus toward $R$.  Arms beginning at a scale comparable
to $K l$, where $K$ is a sufficiently large fixed number, can have
separated inward protrusions and can be required to reach a little
beyond the target cut.  Their probability is bounded below by
$c\pi_2(l,R)$.  To check the relative cost of the landing failure,
crop the remaining arms through a fixed-ratio buffer before the
annulus where the landing test is performed.  Collar comparison
and quasimultiplicativity bound the joint probability of that
failure and the remaining arms by
$C\varepsilon\pi_2(l,R)$.  Choose $\varepsilon$ small, or use the
equivalent well-separated-arm comparison.  This also shows that
the construction is uniform under conditioning on data separated
from it by the specified collars.

The two sets of protrusions can be joined through the free space
between them.  Group their locations on a sufficiently fine but
fixed grid along each cut.  For every resulting pair of location
choices, there are disjoint corridors joining the primal tabs and
the dual tabs, with widths bounded below by a positive multiple
of $l$.  The constants can be small functions of $h$; they do not
change the total number of geometric trial scales.  In a
protrusion box, require a transverse crossing at a level between
the landing cut and the protrusion's outermost required level.
It must meet the protruding arm.  Overlapping longitudinal and
transverse crossings then carry that contact through its corridor.
Only one color is involved in each corridor and its overlap boxes.
The conditional crossing estimates and FKG in that color give a
positive fixed gluing probability.  This is the tube construction
of \cite[Proof of Proposition~5.4]{CDCH}.

The distinction between the central band and its extensions is
essential here.  An arbitrary event $E$ may fix \emph{every edge}
of the central band, so no gluing edge is allowed to lie there.
The protrusions extend outside that band; their overlap boxes and
the connecting corridors lie strictly outside it.  Within an
overlap box the only requirement inherited from the extension is
a connection in that box's color, which is increasing in its
color edges.  Thus FKG applies after $E$ is fixed.  The auxiliary
extensions are integrated out when we take a measure on the
central-band configurations.  They are not additional band data
that must agree in the eventual coupling.

\subsection{Frequent opportunities under the conditioned law}

We now prove Proposition~\ref{prop:outward}.  Denote the common
ambient law, including its permitted exterior reweighting, by
$\Pp_0$, and write
\[
 Q_a=\Pp_0(\,\cdot\mid S_R(a))
\]
for a particular start $a$.  Until the target scale, all conditional
edge laws before the $S_R(a)$ conditioning are ordinary FK laws
with induced partitions, or mixtures of such laws.  Consequently
the collar and crossing comparisons just established apply to
$\Pp_0$ as well.

Choose a fixed large ratio $K$ leaving room for the past-extension
annulus, the reference band with its two collars, and the
outer-extension annulus.  Trial radii are
$l_i=K^i l_0$, where $l_0$ is a fixed large multiple of $s$,
increased to a fixed minimum mesh scale when necessary.  Stop the
trials before a fixed fraction of $R$.  The number $k$ of available
trials satisfies
\[
 k\ge \frac{\log(R/s)}{\log K}-C.
\]
When the right side is bounded, enlarging the constant in
\eqref{eq:outward-error} suffices.  We may therefore assume there
are many trials.

Let $\mathcal F_i$ contain all edges inside the radius-$l_i$ cut,
together with the identities of the start-connected color
regions there.  Call trial $i$ good for start $a$ if
\begin{equation}\label{eq:good-trial}
 \Pp_0(S_{8l_i}(a)\mid\mathcal F_i)\ge h.
\end{equation}
The threshold $h>0$ is chosen after $K$.  The preceding separation
and gluing constructions then use narrower corridors if necessary;
they require no increase of $K$ depending on $h$.

Under $Q_a$, with probability at least $1-Ce^{-c k}$, at least
$3k/4$ of the trials are good.  Here is the conditioning calculation.
The event $S_R(a)$ implies $S_{8l_i}(a)$ at every trial.  For a
specified collection $J$ of bad trials, expose success to
$8l_i$ before passing to the next selected trial.  Since badness is
$\mathcal F_i$-measurable, the tower property and the negation of
\eqref{eq:good-trial} give a factor at most $h$ for each $i\in J$.
Moreover, by \eqref{eq:start-probabilities} and
quasimultiplicativity,
\[
 \frac{\Pp_0(S_R(a))}{\Pp_0(S_{l_0}(a))}
       \ge C^{-1}e^{-C_0 k}.
\]
The microscopic or one-arm factors in
\eqref{eq:start-probabilities} cancel in this ratio.  Summing over
collections of at least $k/4$ bad trials gives
\begin{equation}\label{eq:good-trial-tail}
 Q_a\{\text{at least }k/4\text{ bad trials}\}
 \le C e^{C_0 k}2^k h^{k/4}.
\end{equation}
Choose $h$ so small that the right side decreases exponentially
in $k$.  This choice is uniform over all the permitted starts and
exterior laws.

\subsection{A common band measure and loss of the past}

At a good trial, place a scaled copy of the reference band between
the past-extension and outer-extension annuli.  Its central
configuration has a subprobability measure $\nu_l$ defined as
follows: take the plane FK law, restrict to the successful
reference-band construction of Lemma~\ref{lem:reference-band},
including the auxiliary protrusions, and then retain only the
central-band edges.  There is a fixed $v_0>0$ such that
$\nu_l(\text{all band data})\ge v_0$.

For every event $E$ of central-band data and every good past,
\begin{equation}\label{eq:band-minorization-raw}
 \Pp_0\{\text{band data}\in E,\ S_R(a)
                     \mid\mathcal F_i\}
       \ge c_h\,\nu_l(E)\,\pi_2(l,R).
\end{equation}
Indeed, the good past supplies the separated start-connected
extensions with probability bounded below in terms of $h$.
Require the reference-band event with its protrusions, and
separated arms continuing toward $R$.  Before gluing, these
observations are separated by full collars.  Bounded-density
comparison can therefore be applied successively from the two
sides of the central band.  It preserves the factor $\nu_l(E)$
for an arbitrary, possibly nonmonotone, set $E$, and gives the
factor $\pi_2(l,R)$ for the remaining arms.  There is still
unobserved space for the two-color tube construction.  Its
same-color overlap requirements leave $E$ fixed, as explained
above, and complete both paths from the start to the target.
This proves \eqref{eq:band-minorization-raw}.

To make the averaging explicit, condition beyond a collar
containing all these constructions, and let $\xi$ denote the
resulting exterior state.  Given the already fixed past, write
$p_\xi$ for the conditional probability of reaching $8l$ and
$\mu$ for the conditional distribution of $\xi$.  Goodness
means $\int p_\xi\,\mu(\dd\xi)\ge h$; it does not assert
$p_\xi\ge h$ for each $\xi$.  Uniform collar comparison makes
the landing-error probability at most $\varepsilon$ for every
$\xi$, with $\varepsilon<h/2$.  In the preceding construction,
keep the probability of the start-connected extensions as a
factor instead of replacing it by a constant.  The finitely
many landing-location choices all have the same lower gluing
bound.  The conditional lower bound is therefore
\[
 c\nu_l(E)\pi_2(l,R)(p_\xi-\varepsilon)_+.
\]
Its integral is at least
$c\nu_l(E)\pi_2(l,R)(h-\varepsilon)$, as required in
\eqref{eq:band-minorization-raw}.  This proves the assertion for
the exterior mixture without imposing pointwise goodness.
No identification jumps an annulus used in these comparisons.

In the reverse direction, any successful continuation from the
past must have two arms in the remaining annulus.  Cropping
through a collar and using \eqref{eq:start-probabilities} gives
\begin{equation}\label{eq:band-denominator}
 \Pp_0(S_R(a)\mid\mathcal F_i)\le C\pi_2(l,R).
\end{equation}
Dividing \eqref{eq:band-minorization-raw} by
\eqref{eq:band-denominator} shows that, under $Q_a$, the band
marginal after every good past dominates $c\nu_l$.  Hence two
good pasts, even for different starts, have a uniformly positive
common part on identical central-band data.

It remains to check that coupling this band is enough.  Cut its
strands at both rims.  There are two through-arcs.  All remaining
arcs have both ends on the same rim.  Whatever the configuration
inside the inner rim, completing its pairings joins the two
through-arcs to one another: they are the only two strand ends
left after the same-rim arcs and the interior pairings are
followed.  Interior loops can change in number, but their loop
weight is then a factor depending solely on that interior
completion.  The effective pairing seen on the outer rim is
fixed by the band.

The same fact can be expressed in FK terms.  Exactly one primal
region and one dual region transmit through the band.  An
exterior connection that enters the past must do so through the
transmitted region of its color, whose trace on the outer rim is
already determined by the band.  Thus the induced exterior
primal partition is independent of the earlier past.  On $S_R(a)$,
the starting regions have made contact with these two transmitted
regions.  Conditional on that contact, the remaining constraint
is only that the transmitted regions continue to the target
cut.  Paths that cross the band repeatedly use the same
transmitted region and impose no additional constraint.  The
domain Markov property now shows that the conditional law on
and outside the outer rim, given the common band and successful
continuation, is identical for the two pasts.  A common exterior
weight or an ambient wall farther out changes this common law
in the same way in both cases.

We can therefore complete the coupling as follows.  At each trial,
if both pasts are good, use the common submeasure $c\nu_l$ to
couple the central-band configurations with probability at least
$p>0$.  On success, fill the interior gaps according to their
respective conditional laws and sample the exterior from the
common conditional law just described.  On failure, use the
residual band marginals and fill the data up to the next trial
cut with their correct conditional laws.  No edges beyond that
next cut are inspected merely to decide whether a coupling
attempt failed.  This is a coupling of the successive conditional
kernels, so each complete marginal remains $Q_a$ throughout.

For completeness, this iteration does not require the trials to
be independent.  By \eqref{eq:good-trial-tail} and a union bound
over the two marginals, except with probability $Ce^{-c k}$
there are at least $k/2$ jointly good trials.  Their goodness is
known from the two pasts before the corresponding attempt.
At each such trial, conditional on the entire coupled past and
on all earlier failures, the overlap probability is at least
$p$.  Iterating at these predictable trial times bounds the
probability of $k/2$ failed good attempts by $(1-p)^{k/2}$.
The total failure probability is consequently at most
$Ce^{-c'k}\le C(s/R)^c$.  On success the band and its exterior
agree, and the two transmitted color regions have been reached
by both starts.  This proves Proposition~\ref{prop:outward}.

For a tile-side start, the common through-strand passes through
the associated medial segment: otherwise that segment would lie
on a completed inner loop surrounding one of its two end sites,
contrary to $S_R(e)$.  For a disk start, the common through-strand
meets the disk, since every disk-to-target strand must pass
through the band.  Thus an exterior rule selecting the
distinguished Dobrushin strand gives the same verdict for a
microscopic segment and for a disk after successful coupling.
This is the form used in the occupation comparison.
